\chapter{DP Patterns}\label{ch:dppatterns}

Most dynamic programs belong to a small number of families, and recognising the family gives the state almost for free. A set of
up to 20 items becomes a bitmask; a contiguous segment becomes a pair of endpoints; a huge number becomes its digits plus a
``still equal to the bound'' flag; a grid that is too wide to enumerate becomes the boundary between its filled and empty parts;
a recurrence applied $10^{18}$ times becomes a matrix power. This chapter goes through these families, each with a problem solved
end to end, together with the two topics that most often hide inside them: expectations and games.

\begin{chapterintro}{\gold\ \plat}{Chapter~\ref{ch:dp} (the five-question recipe) and Chapter~\ref{ch:bits} (bitmasks).}
  \item use bitmask DP over subsets, including pair-valued states and submask enumeration in $\Oh{3^n}$;
  \item write interval DP, and turn two-player games into one difference;
  \item count numbers up to $10^{18}$ with digit DP, and tile grids with profile DP;
  \item sum over subsets in $\Oh{n 2^n}$, and convolve over disjoint subsets;
  \item schedule weighted intervals, compute probabilities and expectations, and solve games by DP;
  \item jump $10^{18}$ steps of a linear recurrence with a matrix power.
\end{chapterintro}

% ======================================================================================================================
\section[Bitmask DP]{Bitmask DP \hfill \gold}\label{sec:bitmaskdp}

When the future depends on \emph{which} of $n$ items have been used, and $n$ is at most about 20, store that set as an $n$-bit mask
(Section~\ref{sec:bittricks}). There are $2^n$ sets, about $10^6$ for $n = 20$, so the state may add one more index of size $n$, and
each transition may try every item: $\Oh{2^n n^2}$ is about $4 \cdot 10^8$ simple operations at $n = 20$, which is the usual
budget.

\begin{problem}{Hamiltonian Flights (CSES)}
There are $n \le 20$ cities and $m$ one-way flights. Count the routes from city 1 to city $n$ that visit every city exactly once,
modulo $10^9 + 7$.
\solution $dp[\mathit{mask}][v]$ is the number of routes that start at city 1, visit exactly the cities of $\mathit{mask}$ and end
at $v$. A route grows by one flight $v \to w$ to a city $w \notin \mathit{mask}$. Masks only grow, so increasing numeric order is a
valid order (Figure~\ref{fig:hamilton}). Two observations halve the work and more: every useful mask contains city 1, and city $n$
must be last, so a mask that contains city $n$ but not every city is never extended.
\end{problem}

\begin{figure}[htbp]
\centering
\begin{tikzpicture}
  \tikzset{mk/.style={draw=fgblue, fill=fgsky!14, rounded corners=2pt, font=\ttfamily\small, minimum width=12mm,
    minimum height=6mm, line width=0.7pt}}
  \foreach \nm/\px/\py in {0001/0/0, 0011/2.6/1.2, 0101/2.6/0, 1001/2.6/-1.2, 0111/5.2/1.2, 1011/5.2/0, 1101/5.2/-1.2,
      1111/7.8/0}
    \node[mk] (M\nm) at (\px, \py) {\nm};
  \foreach \nm in {1001, 1011, 1101} \node[mk, dim] at (M\nm) {\nm};
  \node[mk, fill=fgorange!20, draw=fgorange!85!black] at (M1111) {1111};
  \foreach \ua/\ub in {0001/0011, 0001/0101, 0011/0111, 0101/0111, 0111/1111} \draw[ded] (M\ua) -- (M\ub);
  \foreach \ua/\ub in {0001/1001, 0011/1011, 0101/1101} \draw[ded, dashed, draw=fgline] (M\ua) -- (M\ub);
  \node[notel, anchor=west, text width=4.6cm] at (9.0, 0) {Masks for $n = 4$, city 1 being the lowest bit. Every route starts
    at city 1, so only masks with that bit appear; masks with city 4 (the highest bit) that are not complete (grey) receive
    values but are never extended.};
\end{tikzpicture}
\caption{The DP of Hamiltonian Flights moves through subsets of cities, one city per step.}
\label{fig:hamilton}
\end{figure}

\begin{code}
void solve() {
    int n, m;
    cin >> n >> m;
    vvi cnt(n, vi(n, 0));  // the number of flights from city a to city b
    rep(i, m) {
        int a, b;
        cin >> a >> b;
        cnt[a - 1][b - 1]++;
    }
    // dp[mask * n + v]: routes from city 0 through exactly the cities of
    // mask, ending at v; 32-bit entries, since there are 2^20 * 20 of them
    vector<uint32_t> dp((1 << n) * n, 0);
    dp[1 * n + 0] = 1;  // the route that has only visited city 0
    int full = (1 << n) - 1;
    rep(mask, 1, 1 << n) {
        if (!(mask & 1)) continue;  // every route starts at city 0
        if (mask >> (n - 1) & 1 && mask != full) continue;  // n - 1 is last
        rep(u, n) {
            ll ways = dp[mask * n + u];
            if (!ways) continue;
            rep(v, n) if (!(mask >> v & 1) && cnt[u][v]) {  // push to v
                uint32_t& t = dp[(mask | 1 << v) * n + v];
                t = (t + ways * cnt[u][v]) % MOD;
            }
        }
    }
    print(dp[full * n + n - 1]);
}
\end{code}

The table has $2^{20} \cdot 20 \approx 2 \cdot 10^7$ entries: 84\,MB with 32-bit entries, twice that with 64-bit ones, which is why
the type is chosen explicitly. With all 380 possible flights the program runs in 0.18\,s. A flat array indexed by $\mathit{mask}
\cdot n + v$ is also much faster than a \cc{vector} of $2^{20}$ small vectors.

\usetemplate{Bitmask DP (TSP)}{Ctrl+Alt+T, ]}{cptsp}

\subsection{Pair-valued states}

The state does not have to hold a single number. When a set of items is packed into bins one item at a time, the only thing the
future needs is the number of bins used and how full the current bin is.

\begin{problem}{Elevator Rides (CSES)}
There are $n \le 20$ people with weights $w_i$ and an elevator with capacity $x$. What is the fewest rides that carry everyone?
\solution For every set of people, keep the best pair $(\text{rides}, \text{weight of the last ride})$ that carries exactly that
set, compared lexicographically. Adding person $p$ either fits into the last ride or opens a new one. Fewer rides is always better,
and with as many rides a lighter last ride is better, because anything that can board after the heavier can board after the
lighter: so the lexicographic minimum dominates every other way of carrying the set (Table~\ref{tab:elevator}).
\end{problem}

\begin{table}[htbp]
\centering\small
\begin{tabular}{@{}lccc@{}}
\toprule
people already carried & their best pair & boards last & new pair \\
\midrule
$\{4, 8, 6\}$ & $(2, 8)$ & 1 & $(2, 9)$ \\
$\{4, 8, 1\}$ & $(2, 4)$ & 6 & $(2, 10)$ \\
$\{4, 6, 1\}$ & $(2, 1)$ & 8 & $(2, 9)$ \\
$\{8, 6, 1\}$ & $(2, 6)$ & 4 & $(2, 10)$ \\
\midrule
\multicolumn{3}{@{}l}{all four people, the minimum:} & $\mathbf{(2, 9)}$ \\
\bottomrule
\end{tabular}
\caption{Elevator Rides with weights 4, 8, 6, 1 and capacity 10: the state of the full set, from its four predecessors. Two rides
suffice, the last one carrying 9.}
\label{tab:elevator}
\end{table}

\begin{code}
void solve() {
    int n, x;
    cin >> n >> x;
    vi w(n);
    for (ll& v : w) cin >> v;
    // best[mask]: the smallest (rides, weight of the last ride) for carrying
    // exactly the people in mask; the last ride may still take more people
    vector<pll> best(1 << n, {INF, 0});
    best[0] = {1, 0};
    rep(mask, 1, 1 << n) rep(p, n) {
        if (!(mask >> p & 1)) continue;
        auto [rides, last] = best[mask ^ (1 << p)];  // p is the last to board
        if (last + w[p] <= x) last += w[p];
        else rides++, last = w[p];
        chmin(best[mask], pll{rides, last});
    }
    print(best[(1 << n) - 1].first);
}
\end{code}

The obvious alternative, ``$dp[\mathit{mask}]$ is the fewest rides, built from a submask that fits into one ride'', is correct but
costs $\Oh{3^n}$; the pair-valued state brings it to $\Oh{2^n n}$, 0.08\,s for $n = 20$.

\subsection{Enumerating submasks}

Some transitions really do range over all subsets of the current mask, for example ``split the set into groups, each with a
cost''. The loop
\begin{code}
for (int sub = mask; sub > 0; sub = (sub - 1) & mask) { /* use sub */ }
\end{code}
visits every non-empty submask of \cc{mask} in decreasing order: subtracting 1 clears the lowest set bit and sets all bits below
it, and the \cc{&} keeps only those inside the mask. Over all masks the total is $3^n$, not $4^n$, because each element is either
outside the mask, in the mask but not the submask, or in both: three choices per element. That is $3.5 \cdot 10^9$ for $n = 20$,
too many, but fine for $n \le 16$ ($4.3 \cdot 10^7$).

\begin{recognize}
$n \le 20$ (or $n \le 16$ for $\Oh{3^n}$) and an answer that depends on a \emph{set} of chosen or visited things: Hamiltonian
paths, assigning $n$ tasks to $n$ people ($dp[\mathit{mask}]$ is the best cost of giving the first $|\mathit{mask}|$ people the tasks in
$\mathit{mask}$), covering, packing into bins, partitioning into groups.
\end{recognize}

% ======================================================================================================================
\section[Interval DP]{Interval DP \hfill \gold}\label{sec:intervaldp}

When a problem removes, merges or splits \emph{contiguous} parts of a sequence, the state is a segment: $dp[l][r]$ is the answer for
$a_l, \dots, a_r$. The transition either removes one end ($dp[l+1][r]$ and $dp[l][r-1]$) or tries every split point $k$,
\[ dp[l][r] = \min_{l \le k < r} \bigl(dp[l][k] + dp[k+1][r]\bigr) + \mathit{cost}(l, r), \]
and the states are computed by increasing length, so that every shorter segment is ready. Splits cost $\Oh{n^3}$ in total, and
Knuth's optimization (Section~\ref{sec:knuth}) often brings them to $\Oh{n^2}$.

\begin{problem}{Removal Game (CSES)}
Two players take turns removing the first or the last number of a list ($n \le 5000$); each wants to maximise the sum of their own
numbers. What is the first player's total when both play optimally?
\solution Store the \emph{difference} $d[l][r]$, the best (mover's total minus the other's total) on $a_l, \dots, a_r$. The mover
takes an end and then becomes the other player of the rest: $d[l][r] = \max(a_l - d[l+1][r],\ a_r - d[l][r-1])$
(Figure~\ref{fig:removal}). The first player's total is $(\text{sum} + d[0][n-1])/2$.
\end{problem}

\begin{figure}[htbp]
\centering
\begin{tikzpicture}
  \tikzset{tc/.style={draw=fgline, line width=0.5pt, minimum size=8mm, inner sep=0pt, font=\ttfamily\small, anchor=center}}
  \foreach \val [count=\cc from 0] in {4, 5, 1, 3}
    \node[font=\sffamily\small\bfseries, text=fgblue] at (0.8*\cc, 0.75) {\val};
  \node[lbl, anchor=east] at (-0.5, 0.75) {$a$};
  \foreach \rr/\row in {0/{4,1,0,3}, 1/{5,4,3}, 2/{1,2}, 3/{3}}
    \foreach \val [count=\cc from 0] in \row {
      \pgfmathtruncatemacro{\col}{\cc + \rr}
      \node[tc] (D\rr\col) at (0.8*\col, -0.8*\rr) {\val};
    }
  \foreach \rr in {0, 1, 2, 3} \node[lbl, anchor=east] at (-0.5, -0.8*\rr) {$l = \rr$};
  \node[tc, hot] at (D03) {3};
  \node[tc, cool] at (D13) {3};
  \node[tc, cool] at (D02) {0};
  \draw[->, line width=1.1pt, draw=fgverm, shorten <=2.3mm, shorten >=2.3mm] (D13.center) -- (D03.center);
  \draw[->, line width=1.1pt, draw=fgverm, shorten <=2.3mm, shorten >=2.3mm] (D02.center) -- (D03.center);
  \node[notel, anchor=west, text width=6.4cm] at (3.7, -1.2) {$d[l][r]$ for $a = 4, 5, 1, 3$; row $l$, column $r$. For the
    whole list: taking 4 leaves $d[1][3] = 3$ to the opponent, $4 - 3 = 1$; taking 3 leaves $d[0][2] = 0$, $3 - 0 = 3$. So
    $d[0][3] = 3$, and the first player collects $(13 + 3)/2 = 8$.};
\end{tikzpicture}
\caption{The difference table of Removal Game; each cell looks at the cell below it and the cell to its left.}
\label{fig:removal}
\end{figure}

\begin{code}
void solve() {
    int n;
    cin >> n;
    vi a(n);
    ll total = 0;
    for (ll& x : a) cin >> x, total += x;
    // d[l]: the best (mover's total - other's total) on the segment of the
    // current length that starts at l; lengths grow, so d[l + 1] and d[l]
    // still hold the values for the segments one shorter
    vi d(a);
    rep(len, 2, n + 1) rep(l, n - len + 1)
        d[l] = max(a[l] - d[l + 1], a[l + len - 1] - d[l]);
    print((total + d[0]) / 2);
}
\end{code}

The full table would need $n^2 = 2.5 \cdot 10^7$ numbers; going by length, one array suffices, because the segment of length
$\ell$ at $l$ needs only the two segments of length $\ell - 1$ at $l$ and $l + 1$, and $l + 1$ has not been overwritten yet when
$l$ is computed. The program runs in 0.07\,s.

\usetemplate{Range DP}{Ctrl+Alt+T, -}{cprangedp}

\begin{trick}[One number for two players]
In a two-player zero-sum game, store the \emph{difference} between the player to move and the other player. Since the roles swap
after every move, the recurrence is always ``my gain now minus the opponent's best difference afterwards'', and a single table
replaces two. The same trick turns any ``both play optimally'' DP into a maximisation.
\end{trick}

\begin{recognize}
Merging adjacent things (piles of stones, matrix chain multiplication), removing elements whose neighbours then change, palindromes
(``fewest insertions to make a palindrome''), parenthesisations, and games that take from either end. $n \le 500$ is typical for
$\Oh{n^3}$, $n \le 5000$ for $\Oh{n^2}$.
\end{recognize}

% ======================================================================================================================
\section[Digit DP]{Digit DP \hfill \plat}\label{sec:digitdp}

``How many integers in $[a, b]$ have property $P$?'' with $b$ up to $10^{18}$ is answered as $f(b) - f(a - 1)$, where $f(x)$
counts the good integers in $[0, x]$. To count up to $x$, build the integer digit by digit from the most significant one. As long as
every digit so far equals $x$'s digit, the integer is \emph{tight}: the next digit may not exceed $x$'s next digit. The first time a
smaller digit is chosen, the integer is already below $x$, every later digit is free, and the count of completions no longer depends
on $x$ at all (Figure~\ref{fig:digitdp}). So the state is (position, tight, and whatever the property needs, such as the previous
digit), only a few hundred states for 18 digits.

\begin{figure}[htbp]
\centering
\begin{tikzpicture}
  \tikzset{dg/.style={draw=fgverm, fill=fgverm!12, line width=0.9pt, minimum size=7mm, font=\ttfamily\small},
    br/.style={draw=fgblue, fill=fgsky!14, rounded corners=2pt, line width=0.6pt, font=\sffamily\scriptsize, align=center,
      text width=2.55cm, inner sep=3pt}}
  \foreach \dd [count=\ci] in {3, 5, 2, 7} \node[dg] (T\ci) at ({2.9*(\ci - 1)}, 0) {\dd};
  \foreach \ci in {1, 2, 3} {
    \pgfmathtruncatemacro{\nx}{\ci + 1}
    \draw[->, line width=1.2pt, draw=fgverm] (T\ci) -- (T\nx);
  }
  \foreach \tx [count=\ci] in {{first digit 1 or 2,\\then 3 free digits:\\$2 \cdot 9^3 = 1458$},
      {after 3: 0, 1, 2 or 4,\\then 2 free digits:\\$4 \cdot 9^2 = 324$},
      {after 35: 0 or 1,\\then 1 free digit:\\$2 \cdot 9 = 18$},
      {after 352: six digits\\below 7, not 2:\\$6$}}
    { \node[br, anchor=north] (B\ci) at ({2.9*(\ci - 1)}, -1.0) {\tx};
      \draw[->, line width=0.8pt, draw=fgblue, dashed] (T\ci.south) -- (B\ci.north); }
  \node[lbl, anchor=west] at (9.2, 0.0) {$+1$: 3527 itself};
  \node[notel, anchor=west, text width=13cm] at (-0.4, -3.05) {Shorter numbers add $9 + 81 + 729 = 819$ and the number 0 adds 1, so
    $f(3527) = 1 + 819 + 1458 + 324 + 18 + 6 + 1 = 2627$ numbers in $[0, 3527]$ have no two equal adjacent digits.};
\end{tikzpicture}
\caption{Digit DP for Counting Numbers with $x = 3527$. The tight path (red) follows $x$; each dashed branch leaves it with a
smaller digit, after which all later digits are free, so its count is a closed formula.}
\label{fig:digitdp}
\end{figure}

\begin{problem}{Counting Numbers (CSES)}
Count the integers in $[a, b]$ ($0 \le a \le b \le 10^{18}$) in which no two adjacent digits are equal.
\solution For this property the free completions have a closed form: after a fixed digit, each later digit avoids only its left
neighbour, so $k$ free digits give $9^k$ completions. The count is then a walk along the tight path with one branch per position,
exactly the sum of Figure~\ref{fig:digitdp}. For properties without a closed form, the same walk with a memoized DP over (position,
tight, extra state) is the general method, and the template implements it.
\end{problem}

\begin{code}
void solve() {
    ll a, b;
    cin >> a >> b;
    // the numbers in [0, x] with no two equal adjacent digits
    auto count = [](ll x) -> ll {
        if (x < 10) return x + 1;  // also right for x = -1: none
        string s = to_string(x);
        int L = sz(s);
        vi pw(L + 1, 1);  // 9^k: each later digit avoids only its neighbour
        rep(i, 1, L + 1) pw[i] = pw[i - 1] * 9;
        ll res = 1;  // the number 0
        rep(len, 1, L) res += pw[len];  // every shorter length
        rep(i, L) {  // length L, first smaller than x at digit i
            int lo = i == 0 ? 1 : 0;
            rep(d, lo, s[i] - '0')
                if (i == 0 || d != s[i - 1] - '0') res += pw[L - 1 - i];
            if (i > 0 && s[i] == s[i - 1]) return res;  // x's prefix is bad
        }
        return res + 1;  // x itself
    };
    print(count(b) - count(a - 1));
}
\end{code}

\usetemplate{Digit DP}{Ctrl+Alt+T, =}{cpdigit}

\begin{pitfall}[Leading zeros and the empty prefix]
When the property depends on the digits actually written (``no two equal adjacent digits'', ``the digit sum is prime''), leading
zeros must not count as digits: $0012$ is the number 12, whose first digit is 1. Either count each shorter length separately, as
above, or add a \texttt{started} flag to the state. Also check $f(a - 1)$ for $a = 0$: it must be 0, not an error.
\end{pitfall}

% ======================================================================================================================
\section[Boundaries and profiles]{Boundaries and profiles \hfill \plat}\label{sec:profiledp}

When a grid is filled piece by piece, the part already filled matters to the future only through its \emph{boundary}. If the grid is
narrow, the boundary is small enough to be the state.

\paragraph{A two-state boundary.} In Counting Towers (CSES), a tower of width 2 and height $n$ is cut into rectangular blocks; count
the ways. Build it row by row from the bottom. What the next row may do depends only on whether the top row is one block of width 2
(state $a$) or two separate blocks of width 1 (state $b$), as in Figure~\ref{fig:towers}. From $a$, the next row either continues
the wide block, starts a new wide block, or starts two narrow ones; from $b$, each column continues its block or starts a new one
($2 \cdot 2 = 4$ ways), or both blocks end and a wide one starts. So $a_{h+1} = 2a_h + b_h$ and $b_{h+1} = a_h + 4b_h$, with
$a_1 = b_1 = 1$.

\begin{figure}[htbp]
\centering
\begin{tikzpicture}
  \tikzset{blk/.style={draw=fgink!80, line width=0.7pt, fill=#1}, blk/.default=fgsky!25}
  \draw[blk=fgsky!30] (0, 0) rectangle (1.2, 0.55);
  \draw[blk=fgsky!12] (0, -0.55) rectangle (0.6, 0);
  \draw[blk=fgsky!12] (0.6, -0.55) rectangle (1.2, 0);
  \node[blbl] at (0.6, -0.95) {state $a$};
  \node[lbl] at (0.6, -1.3) {top: one wide block};
  \draw[blk=fgorange!30] (5.0, 0) rectangle (5.6, 0.55);
  \draw[blk=fgorange!30] (5.6, 0) rectangle (6.2, 0.55);
  \draw[blk=fgorange!12] (5.0, -0.55) rectangle (6.2, 0);
  \node[blbl] at (5.6, -0.95) {state $b$};
  \node[lbl] at (5.6, -1.3) {top: two narrow blocks};
  \draw[->, line width=1pt, draw=fgblue] (1.5, 0.35) to[bend left=25] node[above, font=\sffamily\scriptsize] {1 way} (4.7, 0.35);
  \draw[->, line width=1pt, draw=fgorange!85!black] (4.7, -0.2) to[bend left=25]
    node[below, font=\sffamily\scriptsize] {1 way} (1.5, -0.2);
  \draw[->, line width=1pt, draw=fgblue] (-0.1, 0.45) arc (60:300:0.45);
  \node[font=\sffamily\scriptsize, anchor=east] at (-0.85, 0.06) {2 ways};
  \draw[->, line width=1pt, draw=fgorange!85!black] (6.3, -0.45) arc (240:480:0.45);
  \node[font=\sffamily\scriptsize, anchor=west] at (7.05, -0.06) {4 ways};
  \node[notel, anchor=west, text width=4.3cm] at (8.2, -0.2) {$a_{h+1} = 2a_h + b_h$\\$b_{h+1} = a_h + 4b_h$\\[2pt]
    Height 2: $3 + 5 = 8$ towers.};
\end{tikzpicture}
\caption{Counting Towers as a two-state machine: the future of a partial tower depends only on the shape of its top row.}
\label{fig:towers}
\end{figure}

\begin{code}
void solve() {
    int t;
    cin >> t;
    const int N = 1000000;
    // a[h]: towers of height h whose top row is one block of width 2;
    // b[h]: towers whose top row is two different blocks of width 1.
    // From a, the next row grows the wide block or starts a new one (to a),
    // or starts two narrow blocks (to b); from b, each column grows its block
    // or starts a new one (4 ways, to b), or a new wide block starts (to a).
    vi a(N + 1), b(N + 1);
    a[1] = b[1] = 1;
    rep(h, 1, N) {
        a[h + 1] = (2 * a[h] + b[h]) % MOD;
        b[h + 1] = (a[h] + 4 * b[h]) % MOD;
    }
    while (t--) {
        int n;
        cin >> n;
        print((a[n] + b[n]) % MOD);
    }
}
\end{code}

\paragraph{Broken profiles.} For a wider grid the boundary is a bitmask. In Counting Tilings (CSES), an $n \times m$ grid with
$n \le 10$ is tiled with dominoes. Decide the cells one at a time, column by column and top to bottom within a column. The
\emph{profile} is the set of $n$ cells just beyond the decided ones: in the current column below the current cell, and in the next
column above it (Figure~\ref{fig:profile}). Bit $r$ of the state says whether that cell is already covered by a domino sticking out
of a decided cell. At the current cell there are three cases: it is already covered (clear the bit and move on), a horizontal domino
covers it and the cell to its right (set the bit), or a vertical domino covers it and the cell below (set the next bit, if that
cell is free). That is $nm\,2^n$ steps of constant work, $10^7$ at the limits.

\begin{figure}[htbp]
\centering
\begin{tikzpicture}
  \tikzset{gq/.style={draw=fgline, line width=0.4pt, minimum size=7mm, inner sep=0pt, anchor=center}}
  \foreach \cc in {0, ..., 5} \foreach \rr in {0, ..., 3} \node[gq] (Q\cc\rr) at (0.7*\cc, -0.7*\rr) {};
  \foreach \cc in {0, 1} \foreach \rr in {0, ..., 3} \node[gq, fill=fgline!45] at (Q\cc\rr) {};
  \foreach \rr in {0, 1} \node[gq, fill=fgline!45] at (Q2\rr) {};
  \foreach \cl in {30, 31, 22, 23} \node[gq, draw=fgblue, line width=1.2pt] at (Q\cl) {};
  \node[gq, hot] at (Q22) {};
  \draw[draw=fgink!80, fill=fgorange!30, rounded corners=2pt, line width=0.6pt] (1.12, 0.26) rectangle (2.38, -0.26);
  \draw[draw=fgink!80, fill=fgorange!30, rounded corners=2pt, line width=0.6pt] (0.42, -1.84) rectangle (1.68, -2.36);
  \node[notel, anchor=west, text width=7.0cm] at (4.6, -1.05) {Grey: decided. Blue frames: the profile, here the cells $(0,
    3)$, $(1, 3)$, $(2, 2)$, $(3, 2)$, one per row. Two dominoes (orange) stick into it, so the state is the mask 1001 (rows 3
    and 0 covered). The current cell (red) is free: place a horizontal domino into column 3, or a vertical one over row 3, which
    is not possible here because that cell is covered.};
\end{tikzpicture}
\caption{The broken profile of Counting Tilings, column 2, row 2 of a $4 \times 6$ grid.}
\label{fig:profile}
\end{figure}

\begin{code}
void solve() {
    int n, m;
    cin >> n >> m;
    // cells are decided column by column, top to bottom; bit r of the mask
    // tells whether cell r of the frontier is already covered (rows above
    // the current cell: in the next column, the others: in this column)
    vi dp(1 << n, 0), nd(1 << n);
    dp[0] = 1;
    rep(c, m) rep(r, n) {
        fill(all(nd), 0);
        rep(mask, 1 << n) {
            ll val = dp[mask];
            if (!val) continue;
            if (mask >> r & 1) {  // covered from the left: move on
                int to = mask ^ 1 << r;
                nd[to] = (nd[to] + val) % MOD;
                continue;
            }
            int to = mask | 1 << r;  // a horizontal domino into column c + 1
            nd[to] = (nd[to] + val) % MOD;
            if (r + 1 < n && !(mask >> (r + 1) & 1)) {  // a vertical domino
                to = mask | 1 << (r + 1);
                nd[to] = (nd[to] + val) % MOD;
            }
        }
        swap(dp, nd);
    }
    print(dp[0]);
}
\end{code}

\usetemplate{Broken profile DP (domino tilings)}{Ctrl+Alt+Y, Shift+P}{cpprofile}

Deciding one \emph{cell} at a time, rather than a whole column at once, is what makes the profile ``broken'': a whole column would
need $2^n \times 2^n$ transitions per column, cell by cell only three per state. The program needs 0.06\,s for a $10 \times 1000$
grid. The same idea counts colourings with local constraints, placements of non-attacking pieces, and paths in narrow grids.

% ======================================================================================================================
\section[Sums over subsets]{Sums over subsets \hfill \plat}\label{sec:sosdp}

Given a value $f(T)$ for every subset $T$ of $B$ bits, the \emph{zeta transform} is $F(S) = \sum_{T \subseteq S} f(T)$ for every
$S$. Enumerating submasks costs $\Oh{3^B}$; the \defn{sum over subsets} (SOS) DP does it in $\Oh{B \cdot 2^B}$ by adding one bit
dimension at a time:
\begin{code}
rep(i, B) rep(m, 1 << B)
    if (m >> i & 1) f[m] += f[m ^ 1 << i];
\end{code}

\emph{Why it works.} After the passes for bits $0, \dots, i$, $f[m]$ holds the sum of the original values of the masks $T$ that agree
with $m$ on bits $i+1, \dots, B-1$ and are submasks of $m$ on bits $0, \dots, i$. The pass for bit $i$ adds, to every $m$ with bit
$i$ set, the value of $m$ without bit $i$, which covers exactly the submasks that have bit $i$ cleared
(Figure~\ref{fig:sos}). After the last pass, $f[m] = F(m)$. Running the passes with $-$ instead of $+$ undoes the transform (the
M\"obius transform), and flipping the condition (\cc{!(m >> i & 1)}, adding $f[m \mid 2^i]$) gives sums over \emph{supersets}.

\begin{figure}[htbp]
\centering
\begin{tikzpicture}
  \tikzset{sm/.style={draw=fgline, line width=0.5pt, fill=white, font=\ttfamily\scriptsize, minimum width=9mm,
    minimum height=4.6mm, inner sep=0pt}}
  \foreach \col/\tx in {0/{$f$}, 1/{after bit 0}, 2/{after bit 1}, 3/{after bit 2: $F$}}
    \node[blbl] at (2.6*\col, 0.55) {\tx};
  \foreach \col in {0, 1, 2, 3} \foreach \mm/\nm in {0/000, 1/001, 2/010, 3/011, 4/100, 5/101, 6/110, 7/111}
    \node[sm] (S\col\mm) at (2.6*\col, -0.6*\mm) {\nm};
  \foreach \col/\bt in {1/1, 2/2, 3/4} {
    \pgfmathtruncatemacro{\prv}{\col - 1}
    \foreach \mm in {0, ..., 7} {
      \draw[line width=0.4pt, draw=fgline] (S\prv\mm.east) -- (S\col\mm.west);
      \pgfmathtruncatemacro{\hasbit}{mod(floor(\mm / \bt), 2)}
      \ifnum\hasbit=1
        \pgfmathtruncatemacro{\src}{\mm - \bt}
        \draw[->, line width=0.8pt, draw=fgblue] (S\prv\src.east) -- (S\col\mm.west);
      \fi
    }
  }
\end{tikzpicture}
\caption{The SOS DP for three bits as a butterfly network. In the pass for bit $i$, every mask with bit $i$ set also receives the
value of the mask without it (blue). After three passes each mask holds the sum over all its submasks; for example 111 has
collected all eight values, along exactly one path each.}
\label{fig:sos}
\end{figure}

\begin{problem}{Bit Problem (CSES)}
Given $n \le 2 \cdot 10^5$ numbers $x_i \le 10^6$, report for every $x_i$ the number of $x_j$ with $x_i \mid x_j = x_i$, the number with
$x_i \mathbin{\&} x_j = x_i$, and the number with $x_i \mathbin{\&} x_j \ne 0$.
\solution The first condition says $x_j$ is a submask of $x_i$: subset sums of the count array. The second says $x_j$ is a
supermask: superset sums. The third is the complement of ``$x_j$ is a submask of $\lnot x_i$''. One pass over 20 bits computes both
tables; 0.34\,s in all.
\end{problem}

\begin{code}
void solve() {
    int n;
    cin >> n;
    const int B = 20;  // x <= 10^6 < 2^20
    vi x(n), sub(1 << B, 0), sup(1 << B, 0);
    for (ll& v : x) cin >> v, sub[v]++, sup[v]++;
    // sub[m]: how many x_j are submasks of m; sup[m]: how many are supermasks
    rep(i, B) rep(m, 1 << B) {
        if (m >> i & 1) sub[m] += sub[m ^ 1 << i];
        else sup[m] += sup[m ^ 1 << i];
    }
    int full = (1 << B) - 1;
    // x_j | x = x, x_j & x = x, and x_j & x != 0 (all but the submasks of ~x)
    for (ll v : x) print(sub[v], sup[v], n - sub[full ^ v]);
}
\end{code}

\usetemplate{SOS DP}{Ctrl+Alt+Y, U}{cpsos}

\subsection{Subset convolution}

The \emph{subset convolution} of $a$ and $b$ is $c(S) = \sum_{T \subseteq S} a(T)\, b(S \setminus T)$: the sum over all ways to split
$S$ into two disjoint parts. Its brute force is $\Oh{3^n}$. The zeta transform alone cannot do it, because the pointwise product of
two zeta transforms counts pairs $(T, U)$ with $T \cup U = S$ that may \emph{overlap}. The fix is to remember sizes: split $a$ by the
number of elements, $a_k(T) = a(T)$ if $|T| = k$ and 0 otherwise, transform every $a_k$ and $b_k$, multiply them as polynomials in
the size, $h_k = \sum_{i + j = k} \hat a_i \hat b_j$, and transform back. A pair $(T, U)$ with $T \cup U = S$ and $|T| + |U| = |S|$
is necessarily disjoint, so $c(S)$ is the inverse transform of $h_{|S|}$ at $S$. The cost is $\Oh{n^2 2^n}$: about 1\,s for $n =
20$ with the template, which also keeps the memory to 176\,MB by storing 32-bit values.

\begin{problem}{Disjoint splits}
Given $a(T)$ and $b(U)$ for every subset of $n \le 20$ elements, compute for every set $S$ the sum of $a(T)\,b(S \setminus T)$ over all
splits of $S$, modulo $10^9 + 7$. (With $a = b = $ the indicator of the ``valid'' groups, $c(S)$ counts the ordered ways to split $S$
into two valid groups.)
\end{problem}

\begin{code}
void solve() {  // cpsubsetconv is pasted above solve()
    int n;
    cin >> n;
    vi a(1 << n), b(1 << n);
    for (ll& x : a) cin >> x;
    for (ll& x : b) cin >> x;
    print(subset_convolution(a, b));
}
\end{code}

\usetemplate{Subset convolution}{Ctrl+Alt+U, Shift+S}{cpsubsetconv}

\begin{recognize}
``For every mask, sum (or count, or take the maximum) over its submasks or supermasks''; pairs with conditions like $x \mathbin{\&} y
= 0$ or $x \mid y = x$; values up to $2^{20}$. Splitting a set into two or more disjoint parts with independent costs is a subset
convolution, and repeating it $k$ times gives the ways to split into $k$ parts (Chapter~\ref{ch:linalg} also covers the XOR and AND
convolutions).
\end{recognize}

% ======================================================================================================================
\section[Scheduling with sorted events]{Scheduling with sorted events \hfill \gold}\label{sec:intervalsched}

When the items are intervals in time, sorting them by \emph{end} time gives a natural order: a solution that ends with interval $i$
consists of $i$ and a solution among the intervals that end before $i$ starts, which form a prefix of the sorted list. Binary search
finds that prefix.

\begin{problem}{Projects (CSES)}
There are $n \le 2 \cdot 10^5$ projects, each with a start day, an end day and a reward. You can work on one project per day, and a
project must be done entirely. Maximise the total reward.
\solution Sort by end day. $dp[i]$ is the best reward using only the first $i$ projects: either project $i$ is skipped, $dp[i-1]$, or
it is taken together with the best of the projects that end before its start day, $dp[j] + p_i$ (Figure~\ref{fig:projects}).
$\Oh{n \log n}$, 0.26\,s at the limits.
\end{problem}

\begin{figure}[htbp]
\centering
\begin{tikzpicture}[x=0.95cm]
  \numline{1}{9}{0}{day}
  \ivl{2}{4}{0.55}{fggreen!65}{$p = 4$}
  \ivl{3}{6}{1.15}{fgline!80}{$p = 6$}
  \ivl{5}{7}{1.75}{fggreen!65}{$p = 3$}
  \ivl{6}{8}{2.35}{fgline!80}{$p = 2$}
  \foreach \yy/\tx in {0.55/{$dp[1] = 4$}, 1.15/{$dp[2] = \max(4, 0 + 6) = 6$}, 1.75/{$dp[3] = \max(6, dp[1] + 3) = 7$},
      2.35/{$dp[4] = \max(7, dp[1] + 2) = 7$}}
    \node[lbl, anchor=west] at (9.6, \yy) {\tx};
  \draw[->, line width=0.8pt, draw=fgblue, dashed] (5, 1.6) to[bend left=30] (4, 0.72);
\end{tikzpicture}
\caption{Projects sorted by end day. The project $[5, 7]$ is compatible with the projects that end before day 5, here only the
first one (dashed), so taking it gives $dp[1] + 3 = 7$. The best schedule takes the two green projects.}
\label{fig:projects}
\end{figure}

\begin{code}
void solve() {
    int n;
    cin >> n;
    vector<array<ll, 3>> pr(n);  // (end day, start day, reward)
    for (auto& p : pr) cin >> p[1] >> p[0] >> p[2];
    sort(all(pr));
    vi ends(n);
    rep(i, n) ends[i] = pr[i][0];
    // dp[i]: the best total reward using only the first i projects by end
    vi dp(n + 1, 0);
    rep(i, n) {
        // projects ending before this one's start day are compatible with it
        int j = lower_bound(all(ends), pr[i][1]) - ends.begin();
        dp[i + 1] = max(dp[i], dp[j] + pr[i][2]);  // skip it, or take it
    }
    print(dp[n]);
}
\end{code}

% ======================================================================================================================
\section[Probability and expectation]{Probability and expectation \hfill \plat}\label{sec:probdp}

A probability DP replaces ``number of ways'' with ``probability'', and the transitions multiply by the probability of each choice.
Expected values add one more idea.

\begin{fact}[Linearity of expectation]
$E[X + Y] = E[X] + E[Y]$ for \emph{any} random variables, independent or not. So the expected number of events that happen is the sum
of their probabilities: the expected number of fixed points of a random permutation is $n \cdot \frac1n = 1$, whatever $n$ is.
\end{fact}

\begin{problem}{Dice Probability (CSES)}
Throw $n \le 100$ dice. What is the probability that the sum is between $a$ and $b$?
\solution The distribution of the sum after $i$ dice is a DP over the sum: each value spreads $\frac16$ of its probability to the six
next sums (Figure~\ref{fig:diceprob}).
\end{problem}

\begin{figure}[htbp]
\centering
\begin{tikzpicture}
\begin{axis}[cpplot, width=12cm, height=5.2cm, xlabel={sum $/$ number of dice}, ylabel={probability $\times$ dice},
  xmin=1, xmax=6, ymin=0, ymax=0.82, legend pos=north east, legend columns=2]
  \addplot+[mark=none] coordinates {(1,0.1667) (2,0.1667) (3,0.1667) (4,0.1667) (5,0.1667) (6,0.1667)};
  \addplot+[mark=none] coordinates {(1,0.0556) (1.5,0.1111) (2,0.1667) (2.5,0.2222) (3,0.2778) (3.5,0.3333) (4,0.2778)
    (4.5,0.2222) (5,0.1667) (5.5,0.1111) (6,0.0556)};
  \addplot+[mark=none] coordinates {(1,0.0139) (1.333,0.0417) (1.667,0.0833) (2,0.1389) (2.333,0.2083) (2.667,0.2917)
    (3,0.3472) (3.333,0.375) (3.667,0.375) (4,0.3472) (4.333,0.2917) (4.667,0.2083) (5,0.1389) (5.333,0.0833) (5.667,0.0417)
    (6,0.0139)};
  \addplot+[mark=none] coordinates {(1.0,0.0) (1.6,0.001) (1.8,0.004) (2.0,0.014) (2.2,0.040) (2.4,0.095) (2.6,0.190)
    (2.8,0.326) (3.0,0.485) (3.2,0.629) (3.4,0.715) (3.5,0.727) (3.6,0.715) (3.8,0.629) (4.0,0.485) (4.2,0.326) (4.4,0.190)
    (4.6,0.095) (4.8,0.040) (5.0,0.014) (5.2,0.004) (5.4,0.001) (6.0,0.0)};
  \legend{1 die, 2 dice, 3 dice, 10 dice}
\end{axis}
\end{tikzpicture}
\caption{The distribution of the sum of $n$ dice, rescaled to the average per die. One die is flat, two give a triangle, and ten
are already close to a bell curve centred at 3.5.}
\label{fig:diceprob}
\end{figure}

\begin{code}
void solve() {
    int n, a, b;
    cin >> n >> a >> b;
    vector<double> p(6 * n + 1, 0.0);  // the distribution of the sum so far
    p[0] = 1;
    rep(i, n) {
        vector<double> q(6 * n + 1, 0.0);
        rep(s, 6 * i + 1) if (p[s] > 0)
            rep(d, 1, 7) q[s + d] += p[s] / 6;
        swap(p, q);
    }
    double ans = 0;
    rep(s, a, b + 1) ans += p[s];
    cout << fixed << setprecision(6) << ans << '\n';
}
\end{code}

\paragraph{Counting with $P(\max \le x)$.} The maximum of independent values is awkward directly but easy through its cumulative
distribution: $\max \le x$ exactly when every value is $\le x$. In Candy Lottery (CSES), each of $n$ children gets a uniform number of
candies in $1..k$, and $P(\max \le x) = (x / k)^n$, so $P(\max = x) = (x/k)^n - ((x - 1)/k)^n$ and the expectation is a single sum.

\begin{code}
void solve() {
    int n, k;
    cin >> n >> k;
    // P(max <= x) = (x / k)^n, so P(max = x) is a difference of two powers
    double ans = 0;
    rep(x, 1, k + 1) {
        double le = pow((double)x / k, n), lt = pow((double)(x - 1) / k, n);
        ans += x * (le - lt);
    }
    cout << fixed << setprecision(6) << ans << '\n';
}
\end{code}

\begin{problem}{Moving Robots (CSES)}
Each square of a chessboard has a robot. In each of $k \le 100$ steps, every robot moves to a uniformly random neighbouring square.
What is the expected number of empty squares at the end?
\solution By linearity, the expected number of empty squares is the sum over squares of the probability that the square is empty.
The robots move independently, so that probability is the product, over all robots, of the probability that the robot is not
there; and where one robot ends is a 64-state probability DP over $k$ steps (Figure~\ref{fig:robots}).
\end{problem}

\begin{figure}[htbp]
\centering
\begin{tikzpicture}
  \foreach \cc/\rr/\pc/\val in {0/0/60/0.54, 1/0/33/0.41, 2/0/32/0.40, 3/0/34/0.41, 4/0/34/0.41, 5/0/32/0.40, 6/0/33/0.41,
      7/0/60/0.54, 0/1/33/0.41, 1/1/8/0.29, 2/1/12/0.31, 3/1/12/0.31, 4/1/12/0.31, 5/1/12/0.31, 6/1/8/0.29, 7/1/33/0.41,
      0/2/32/0.40, 1/2/12/0.31, 2/2/12/0.31, 3/2/14/0.32, 4/2/14/0.32, 5/2/12/0.31, 6/2/12/0.31, 7/2/32/0.40, 0/3/34/0.41,
      1/3/12/0.31, 2/3/14/0.32, 3/3/15/0.32, 4/3/15/0.32, 5/3/14/0.32, 6/3/12/0.31, 7/3/34/0.41, 0/4/34/0.41, 1/4/12/0.31,
      2/4/14/0.32, 3/4/15/0.32, 4/4/15/0.32, 5/4/14/0.32, 6/4/12/0.31, 7/4/34/0.41, 0/5/32/0.40, 1/5/12/0.31, 2/5/12/0.31,
      3/5/14/0.32, 4/5/14/0.32, 5/5/12/0.31, 6/5/12/0.31, 7/5/32/0.40, 0/6/33/0.41, 1/6/8/0.29, 2/6/12/0.31, 3/6/12/0.31,
      4/6/12/0.31, 5/6/12/0.31, 6/6/8/0.29, 7/6/33/0.41, 0/7/60/0.54, 1/7/33/0.41, 2/7/32/0.40, 3/7/34/0.41, 4/7/34/0.41,
      5/7/32/0.40, 6/7/33/0.41, 7/7/60/0.54}
    \node[draw=white, line width=0.6pt, fill=fgblue!\pc, minimum size=7.5mm, inner sep=0pt, font=\sffamily\tiny,
      text=fgink] at (0.75*\cc, -0.75*\rr) {\val};
  \node[notel, anchor=west, text width=5.2cm] at (6.1, -2.6) {The probability that each square is empty after $k = 10$ steps.
    Corner squares have only two neighbours, so robots leave them more often than they arrive. The expected number of empty
    squares is the sum of the 64 values, 23.12.};
\end{tikzpicture}
\caption{Moving Robots: a heat map of the per-square probabilities whose sum is the answer.}
\label{fig:robots}
\end{figure}

\begin{code}
void solve() {
    int k;
    cin >> k;
    // the robots move independently, so a cell is empty with probability
    // prod over robots (1 - P(this robot ends there)); add over the cells
    vector<double> empty(64, 1.0);
    rep(start, 64) {
        vector<double> p(64, 0.0);  // where this robot is after each step
        p[start] = 1;
        rep(step, k) {
            vector<double> q(64, 0.0);
            rep(c, 64) {
                if (p[c] == 0) continue;
                int r = c / 8, col = c % 8;
                vi nb;
                if (r > 0) nb.push_back(c - 8);
                if (r < 7) nb.push_back(c + 8);
                if (col > 0) nb.push_back(c - 1);
                if (col < 7) nb.push_back(c + 1);
                for (int d : nb) q[d] += p[c] / sz(nb);
            }
            swap(p, q);
        }
        rep(c, 64) empty[c] *= 1 - p[c];
    }
    double ans = 0;
    for (double e : empty) ans += e;
    cout << fixed << setprecision(6) << ans << '\n';
}
\end{code}

\paragraph{Expected values with loops.} An expected-value DP usually runs backwards from the end: $E[s]$, the expected remaining cost
from state $s$, is the cost of one step plus $\sum_t P(s \to t)\, E[t]$. When a state can return to itself with probability $p$, the
equation contains $E[s]$ on both sides and is solved for it: rolling a die until a 6 appears takes $E = 1 + \frac56 E$ throws, so $E =
6$, and in general an experiment that succeeds with probability $p$ needs $1/p$ tries on average. When states form longer cycles,
the equations form a linear system (Gaussian elimination, Chapter~\ref{ch:linalg}).

\begin{trick}[Expectations modulo a prime]
When an expected value is requested ``as $P \cdot Q^{-1} \bmod 998244353$'', run the whole DP in modular arithmetic: every division by
a probability's denominator becomes a multiplication by a modular inverse (Chapter~\ref{ch:numbertheory}), and the result is exact.
Floating point is only for problems that print decimals, and then \cc{double} with error $\le 10^{-6}$ is enough far more often than
\cc{long double} is needed.
\end{trick}

% ======================================================================================================================
\section[Games]{Games \hfill \gold}\label{sec:gamedp}

In a two-player game with perfect information, no chance and no draws, where the player who cannot move loses, every position is
either \defn{winning} for the player to move or \defn{losing}. A position is winning exactly when \emph{some} move leads to a losing
position, and losing when \emph{every} move (possibly none) leads to a winning one. Positions that come later in play are decided
first, so this is a DP over the DAG of positions, in reverse order of play.

\begin{problem}{Stick Game (CSES)}
A heap has $n \le 10^6$ sticks, and a move removes $p_1$, $p_2$, \dots or $p_k$ of them ($k \le 100$). For every heap size from 1 to $n$,
does the first player win?
\solution $\mathit{win}[0]$ is false (no move), and $\mathit{win}[i]$ is true when some $p_j \le i$ has $\mathit{win}[i - p_j]$ false
(Figure~\ref{fig:stickgame}): $\Oh{nk}$, 0.11\,s.
\end{problem}

\begin{figure}[htbp]
\centering
\begin{tikzpicture}
  \node[blbl, anchor=east] at (-0.2, 0) {position};
  \arrs{W}{0}{0}{L/hot, W/good, L/hot, W/good, W/good, W/good, W/good, L/hot, W/good, L/hot, W/good, W/good, W/good, W/good,
    L/hot}
  \idx{W}{15}
  \foreach \ua/\ub/\dx in {2/1/0.13, 5/1/-0.13, 6/3/0.13, 7/3/-0.13}
    \draw[->, line width=0.8pt, draw=fggreen!70!black] (W-\ua.north) to[bend right=40] ($(W-\ub.north)+(\dx, 0)$);
  \node[notel, anchor=north west, text width=10.4cm] at (0, -0.9) {Moves remove 1, 3 or 4 sticks. Every winning position (W)
    has a move to a losing one (arrows: from 1, 4, 5 and 6 sticks); from 0, 2 and 7 sticks every move reaches a W. The pattern
    repeats with period 7, which is typical: game DPs over a single number often become periodic.};
\end{tikzpicture}
\caption{Winning and losing positions of the Stick Game with moves $\{1, 3, 4\}$.}
\label{fig:stickgame}
\end{figure}

\begin{code}
void solve() {
    int n, k;
    cin >> n >> k;
    vi p(k);
    for (ll& x : p) cin >> x;
    // win[i]: the player to move with i sticks left wins; with 0 sticks
    // there is no move, so the player to move has lost
    vector<bool> win(n + 1, false);
    rep(i, 1, n + 1) for (ll x : p)
        if (x <= i && !win[i - x]) {  // a move to a losing position
            win[i] = true;
            break;
        }
    string s;
    rep(i, 1, n + 1) s += win[i] ? 'W' : 'L';
    print(s);
}
\end{code}

When a position is a \emph{sum} of independent games (several heaps, the player moves in one of them), the win/lose labels of the
parts do not determine the label of the sum; their Grundy numbers do, and Nim is the model (Chapter~\ref{ch:combinatorics}).

% ======================================================================================================================
\section[Matrix exponentiation]{Linear recurrences and matrix powers \hfill \gold}\label{sec:matexp}

When the state is a small vector and every step transforms it by the same linear map, $n$ steps are the $n$-th power of a matrix,
and fast exponentiation computes it with $\Oh{\log n}$ matrix products. For a state of size $k$ that is $\Oh{k^3 \log n}$: for $k = 6$
and $n = 10^{18}$, about 60 products of $6 \times 6$ matrices.

\begin{problem}{Throwing Dice (CSES)}
Count the ordered ways to reach the sum $n$ with dice, modulo $10^9 + 7$, for $n \le 10^{18}$.
\solution Dice Combinations (Section~\ref{sec:dp-recipe}) with a huge $n$. The recurrence $f(s + 1) = f(s) + f(s - 1) + \dots + f(s -
5)$ moves the vector of the last six values by a fixed matrix (Figure~\ref{fig:matpow}); since $f(0) = 1$ and earlier values are 0,
$f(n)$ is the top-left entry of $T^n$. The program takes 0.08\,s.
\end{problem}

\begin{figure}[htbp]
\centering
\small
$\begin{pmatrix} f(s+1) \\ f(s) \\ f(s-1) \\ f(s-2) \\ f(s-3) \\ f(s-4) \end{pmatrix} =
\underbrace{\begin{pmatrix} 1&1&1&1&1&1 \\ 1&0&0&0&0&0 \\ 0&1&0&0&0&0 \\ 0&0&1&0&0&0 \\ 0&0&0&1&0&0 \\ 0&0&0&0&1&0
\end{pmatrix}}_{T}
\begin{pmatrix} f(s) \\ f(s-1) \\ f(s-2) \\ f(s-3) \\ f(s-4) \\ f(s-5) \end{pmatrix}$
\par\bigskip
\begin{tikzpicture}
  \foreach \ex/\sy [count=\ci from 0] in {1/hot, 2/n, 4/hot, 8/hot}
    \node[draw=fgline, line width=0.6pt, rounded corners=2pt, minimum width=10mm, minimum height=7mm, \sy] (P\ci) at
      (1.35*\ci, 0) {$T^{\ex}$};
  \foreach \ci in {1, 2, 3} {
    \pgfmathtruncatemacro{\pv}{\ci - 1}
    \draw[->, line width=0.8pt, draw=fgline] (P\pv) -- (P\ci);
  }
  \node[notel, anchor=west, text width=6.2cm] at (4.75, 0) {Each power is the square of the one before it.
    $n = 13 = 1101_2$, so $T^{13} = T^{8} \cdot T^{4} \cdot T^{1}$ (red).};
\end{tikzpicture}
\caption{Top: one step of the dice recurrence as a matrix product. Bottom: fast exponentiation squares $T$ repeatedly and
multiplies in the powers that match the set bits of $n$.}
\label{fig:matpow}
\end{figure}

\begin{code}
void solve() {
    ll n;
    cin >> n;
    using Mat = array<array<ll, 6>, 6>;
    auto mul = [](const Mat& A, const Mat& B) {
        Mat C{};
        rep(i, 6) rep(k, 6) if (A[i][k]) rep(j, 6)
            C[i][j] = (C[i][j] + A[i][k] * B[k][j]) % MOD;
        return C;
    };
    // the vector (f(s), f(s - 1), ..., f(s - 5)) moves one step by T
    Mat T{}, R{};
    rep(j, 6) T[0][j] = 1;  // f(s + 1) = f(s) + ... + f(s - 5)
    rep(i, 1, 6) T[i][i - 1] = 1;  // the other entries shift down
    rep(i, 6) R[i][i] = 1;
    for (ll e = n; e > 0; e >>= 1, T = mul(T, T))
        if (e & 1) R = mul(R, T);
    print(R[0][0]);  // T^n applied to (f(0), 0, ..., 0) = (1, 0, ..., 0)
}
\end{code}

\usetemplate{Matrix exponentiation}{Ctrl+Alt+T, Shift+X}{cpmatrix}

\begin{trick}[Beyond addition and multiplication]
A matrix power works over any \emph{semiring}. With $(\min, +)$ in place of $(+, \times)$, the $n$-th power of the weighted adjacency
matrix gives the shortest walks with exactly $n$ edges (Graph Paths~II); with ordinary arithmetic it counts walks of length $n$
(Graph Paths~I); with (OR, AND) it gives reachability in exactly $n$ steps. For a recurrence of order $k$ in the thousands, the
$\Oh{k^3 \log n}$ matrix power is too slow, and the polynomial methods of Chapter~\ref{ch:linalg} (Kitamasa, Berlekamp--Massey) take
$\Oh{k^2 \log n}$ or $\Oh{k \log k \log n}$.
\end{trick}

\begin{recognize}
$n$ up to $10^{18}$, a DP whose state is small ($k \le 100$) and whose transition does not depend on the step: linear recurrences,
counting walks of length $n$ in a small graph, DP over positions with a small alphabet, tilings of a $k \times n$ strip.
\end{recognize}

\begin{summary}[Chapter summary]
\begin{itemize}
  \item Bitmask DP handles sets of up to about 20 items in $\Oh{2^n n^2}$; pair-valued states and dominance can replace $\Oh{3^n}$
        submask enumeration, which is itself $3^n$ and not $4^n$.
  \item Interval DP fills $dp[l][r]$ by increasing length; for games, storing the difference of the two players' scores turns the
        game into a single maximisation.
  \item Digit DP counts up to $10^{18}$ with a tight flag: after the first smaller digit, the count of completions is independent of
        the bound.
  \item A narrow grid is filled cell by cell with its boundary (profile) as a bitmask: $\Oh{nm\,2^n}$.
  \item SOS DP sums over subsets in $\Oh{B\,2^B}$; ranking by size turns it into subset convolution in $\Oh{n^2 2^n}$.
  \item Linearity of expectation reduces an expected count to a sum of probabilities; maxima go through $P(\max \le x)$.
  \item Win/lose positions are a DP in reverse order of play; fixed linear transitions jump $10^{18}$ steps with matrix powers.
\end{itemize}
\end{summary}

\begin{practice}
\ppgroup{Sets, intervals, digits and profiles}
\begin{pplist}
  \item Hamiltonian Flights
  \item Elevator Rides
  \item Removal Game
  \item Counting Numbers
  \item Counting Towers
  \item Counting Tilings
  \item Bit Problem
  \item Projects
\end{pplist}
\ppgroup{Probability, games and matrices}
\begin{pplist}
  \item Dice Probability
  \item Candy Lottery
  \item Moving Robots
  \item Stick Game
  \item Throwing Dice
  \item Fibonacci Numbers
  \item Graph Paths I
  \item Graph Paths II
\end{pplist}
\end{practice}
